\section{Направления дальнейших исследований}

Опираясь на полученные в ходе симуляций экспериментальные данные и подтвержденную эффективность нейроэволюционных подходов, был намечен ряд стратегических векторов для дальнейшего развития данного проекта. В частности, планируется масштабирование разработанных методов как в сторону усложнения архитектуры агентов, так и в сторону практического применения в прикладных симуляторах.

\subsection{Разработка гибридных мультимодальных моделей}

Первым направлением развития является преодоление ограничений классического алгоритма NEAT при работе со средами, обладающими высокой размерностью входных данных. В ходе исследования было отмечено, что базовая нейроэволюция испытывает существенные вычислительные трудности при обработке сырых пикселей (например, при получении визуальной информации напрямую с камер агента). Прямая подача изображений на вход NEAT приводит к комбинаторному взрыву возможных топологий и критическому замедлению эволюции.

Для решения этой проблемы мы планируем разработать и протестировать гибридную мультимодальную архитектуру, в которой задачи будут жестко разделены. В роли модуля зрения выступят сверточные нейронные сети. Их задачей станет автоматическое извлечение визуальных признаков и сжатие исходной картинки до компактного информационного вектора.

Затем этот вектор будет передаваться непосредственно в алгоритм NEAT. Апробацию этой связки планируется провести на стандартизированном наборе сред с визуальным входом, например, классических играх Atari.

\subsection{Интеграция алгоритмов в симулятор популяций}

Второе направление представляет собой глобальную практическую цель, послужившую фундаментальной мотивацией для проведения данного дипломного исследования. В настоящее время на базе института ИСИ СО РАН ведется разработка комплексного симулятора поведения популяций виртуальных муравьев.

Внедрение алгоритмов машинного обучения в подобные масштабные проекты напрямую, «вслепую», несет высокие риски нестабильности системы. В связи с этим было критически важно провести предварительный изолированный анализ: проверить сильные и слабые стороны алгоритма NEAT на стандартизированных кинематических моделях (таких как среда непрерывного управления Ant-v4 в физическом движке MuJoCo), настроить функции вознаграждения и изучить поведение агентов вне влияния побочных факторов.

Мы на практике подтвердили, что компактность сформированных топологий позволяет запускать сотни и тысячи независимых агентов одновременно без перегрузки вычислительных мощностей системы и что разреженная архитектура сетей делает агентов устойчивее к изменениям внешней среды и непредсказуемым физическим взаимодействиям.

В качестве финального этапа развития проекта, протестированный и оптимизированный пайплайн обучения NEAT будет интегрирован в архитектуру институтского симулятора. Это позволит перейти от управления одиночными когнитивными агентами к моделированию сложных эмерджентных форм поведения в рамках целых популяций виртуальных насекомых.